\documentclass[11pt,a4paper]{article}

\usepackage{fontspec}
\setmainfont{DejaVu Serif}[
  BoldFont={DejaVu Serif Bold},
  ItalicFont={DejaVu Serif},
  ItalicFeatures={FakeSlant=0.18},
  BoldItalicFont={DejaVu Serif Bold},
  BoldItalicFeatures={FakeSlant=0.18}
]
\setsansfont{DejaVu Sans}
\usepackage{polyglossia}
\setmainlanguage{english}
\usepackage[a4paper,margin=20.5mm]{geometry}
\usepackage{microtype}
\usepackage{setspace}
\setstretch{1.035}
\usepackage{amsmath,amssymb,amsthm,mathtools,bm}
\usepackage{booktabs,longtable,array,tabularx}
\usepackage{enumitem}
\usepackage{xcolor}
\usepackage[unicode,hidelinks]{hyperref}
\usepackage[nameinlink,noabbrev]{cleveref}
\emergencystretch=3em
\allowdisplaybreaks

\hypersetup{
  pdftitle={Foucault Networks on Celestial Bodies as Typed Observation Geometry},
  pdfauthor={Stas, Independent Research Program},
  pdfsubject={Frame-consistent local Foucault dynamics and nested observation maps},
  pdfkeywords={Foucault pendulum, rotating frame, Coriolis force, celestial hierarchy, torus flow, observation geometry}
}

\definecolor{obsblue}{RGB}{37,83,138}
\definecolor{obsred}{RGB}{150,45,45}
\definecolor{lightblue}{RGB}{239,246,253}
\definecolor{lightred}{RGB}{253,242,242}

\theoremstyle{definition}
\newtheorem{definition}{Definition}[section]
\newtheorem{model}[definition]{Model}
\theoremstyle{plain}
\newtheorem{theorem}[definition]{Theorem}
\newtheorem{proposition}[definition]{Proposition}
\newtheorem{corollary}[definition]{Corollary}
\theoremstyle{remark}
\newtheorem{remark}[definition]{Remark}
\newtheorem{warning}[definition]{Boundary}

\newcommand{\R}{\mathbb{R}}
\newcommand{\T}{\mathbb{T}}
\newcommand{\dd}{\mathop{}\!d}
\newcommand{\norm}[1]{\left\lVert #1\right\rVert}
\newcommand{\cl}{\operatorname{cl}}
\newcommand{\SO}{\mathrm{SO}}
\newcommand{\Ocal}{\mathcal O}
\newcommand{\Kcal}{\mathcal K}
\newcommand{\Rcal}{\mathcal R}
\newcolumntype{L}[1]{>{\raggedright\arraybackslash}p{#1}}
\newcolumntype{Y}{>{\raggedright\arraybackslash}X}

\title{\bfseries Foucault Networks on Celestial Bodies\\
\large Frame-consistent local dynamics and nested observation maps}
\author{Stas\\
\small Independent Research Program; ORCID: 0009-0000-2294-705X}
\date{Strict GO Core reference revision v1.1 --- 28 July 2026}

\begin{document}
\maketitle

\begin{center}
\fcolorbox{obsblue}{lightblue}{%
\parbox{0.92\linewidth}{\small
\textbf{Status.} This document supersedes the bilingual working draft
v1.0. It is migrated to GO Core v0.2 and the Frames--Forces--Dissipation
Interface v0.1. The local pendulum is a small-amplitude rotating-frame
model; the celestial hierarchy is a finite kinematic laboratory, not
an ephemeris or a precision geophysical model.}}
\end{center}

\begin{abstract}
We formulate a network of idealized Foucault pendula on a finite
hierarchy of rotating bodies without mixing body, local, inertial, and
observer-frame components. A body edge in the hierarchy is transformed
through its parent's attitude, a local endpoint displacement is
transformed together with the site position, and a moving observer is
represented by a full affine frame map rather than an origin
subtraction alone. In the body-fixed tangent plane the declared
small-amplitude model is
\[
\ddot u+2\beta\dot u+2\Omega_nJ\dot u+\omega_0^2u=0,
\qquad \Omega_n=\Omega^B\cdot n^B .
\]
For \(\beta=0\), the substitution
\(u_{\mathbb C}=e^{-i\Omega_nt}w\) proves that the oscillation plane
precesses at signed rate \(-\Omega_n\) in the chosen orientation. The
Coriolis term performs no work, while damping dissipates
\(P_{\rm diss}=2m\beta\norm{\dot u}^2\ge0\). A nested
quasiperiodic hierarchy is represented as a continuous readout of a
torus flow. Its closure is the image of the subtorus determined by
integer frequency relations. Frequency support is preserved under a
static Euclidean change of observer and a time-origin shift, but not
under an arbitrary time-dependent rotating observer or a lossy camera
projection. The result is a typed observation model with explicit
claim boundaries, not a new theory of celestial mechanics.
\end{abstract}

\tableofcontents

\section{Scope and observation protocol}

The experiment is a finite mathematical laboratory: choose bodies
\(b\), surface sites \(\ell\), local pendula, an observer frame \(O\),
and a sampling/camera protocol \(\lambda\). The typed chain is
\[
X_t\xrightarrow{\Rcal}\widetilde X_t
\xrightarrow{\Ocal_\lambda}Z_t
\xrightarrow{\Kcal_\lambda}\operatorname{Law}(Y_t\mid Z_t)
\xrightarrow{\widehat\theta_\lambda}\widehat\theta_t.
\]
The reduction \(\Rcal\) selects a finite hierarchy and local model,
\(\Ocal_\lambda\) produces ideal framed endpoints or projected pixels,
\(\Kcal_\lambda\) represents sampling and noise, and
\(\widehat\theta_\lambda\) estimates rates or spectral relations.
These maps are not identified.

\begin{warning}[Claim firewall]
The model excludes tidal elasticity, polar motion, atmospheric
loading, suspension anisotropy, finite-amplitude precession, general
relativity, and precision ephemerides. Parameters called ``planetary''
or ``galactic'' are labels unless supplied by a dated external model
with uncertainty.
\end{warning}

\section{Frame-consistent celestial hierarchy}

Fix an inertial reference \(I\). Each body \(b\) has a center
\(R_b^I(t)\) and attitude
\[
A_{I\leftarrow B_b}(t)\in\SO(3),
\]
where \(B_b\) is its body frame. If \(p(b)\) is the parent and the
relative center vector is stored in the parent frame, recursion is
\begin{equation}\label{eq:tree}
R_b^I
=R_{p(b)}^I
+A_{I\leftarrow B_{p(b)}}\,r_{b\mid p(b)}^{B_{p(b)}}.
\end{equation}
Thus every addition in \cref{eq:tree} occurs in \(I\).

At site \(\ell\), let \(r_{b\ell}^{B_b}\) be the body-fixed site point,
\(n_{b\ell}^{B_b}\) the outward unit normal, and
\[
E_{B_b\leftarrow L_{b\ell}}
=\begin{pmatrix}e_1^{B_b}&e_2^{B_b}&n^{B_b}\end{pmatrix}
\in\SO(3)
\]
the local frame. A horizontal pendulum displacement is
\(u^{L_{b\ell}}=(u_1,u_2,0)^{\mathsf T}\).

\begin{definition}[Endpoint point map]
The endpoint in the inertial frame is
\begin{equation}\label{eq:endpoint-I}
p_{b\ell}^I
=R_b^I+
A_{I\leftarrow B_b}
\left(
r_{b\ell}^{B_b}
+E_{B_b\leftarrow L_{b\ell}}u^{L_{b\ell}}
\right).
\end{equation}
\end{definition}

For a moving observer with origin \(c_O^I(t)\) and attitude
\(A_{I\leftarrow O}(t)\), the observer components are
\begin{equation}\label{eq:observer-frame}
p_{b\ell}^{O}
=A_{O\leftarrow I}
\bigl(p_{b\ell}^I-c_O^I\bigr).
\end{equation}
A camera map \(\Pi_\lambda\) then gives
\(z_{b\ell}=\Pi_\lambda(p_{b\ell}^{O})\). In general
\(\Pi_\lambda\) is not invertible.

\begin{remark}[Correction to v1.0]
The expression
\(R_b^I+A_{I\leftarrow B_b}r_{b\ell}^{B_b}+u_{b\ell}^{B_b}-c_O^I\)
mixed the body displacement \(u_{b\ell}^{B_b}\) with inertial
vectors. Equation \cref{eq:endpoint-I} transforms the site and local
displacement before their inertial addition; \cref{eq:observer-frame}
also includes observer attitude.
\end{remark}

\section{Local rotating-frame pendulum}

Let the body angular velocity be \(\Omega^B\) and set
\[
\Omega_n=\Omega^B\cdot n^B .
\]
On the oriented tangent plane define
\[
Jv=n^B\times v,\qquad J^2=-I .
\]
The following is a declared reduced model, not the exact spherical
pendulum.

\begin{model}[Linear Foucault layer]\label{model:foucault}
For pendulum length \(\ell_p>0\),
\(\omega_0^2=g_{\rm eff}/\ell_p\), and damping rate
\(\beta\ge0\), the body-fixed horizontal displacement obeys
\begin{equation}\label{eq:local-foucault}
\ddot u+2\beta\dot u+2\Omega_nJ\dot u+\omega_0^2u=0.
\end{equation}
The model assumes a fixed local vertical, small displacement,
constant \(\Omega_n\), and neglects terms beyond the horizontal
Coriolis projection.
\end{model}

\begin{theorem}[Signed precession in the undamped reduced model]
\label{thm:precession}
Under \cref{model:foucault} with \(\beta=0\), identify the oriented
tangent plane with \(\mathbb C\), so \(J\) acts as multiplication by
\(i\). Every solution has the form
\[
z(t)=e^{-i\Omega_nt}
\left(C_+\cos(\omega_ct)+C_-\sin(\omega_ct)\right),
\qquad
\omega_c=\sqrt{\omega_0^2+\Omega_n^2}.
\]
Consequently a linearly polarized oscillation plane precesses in the
body frame at signed rate
\[
\boxed{\dot\alpha=-\Omega_n}
\]
for the stated orientation convention.
\end{theorem}

\begin{proof}
With \(z=e^{-i\Omega_nt}w\), the complex form
\[
\ddot z+2i\Omega_n\dot z+\omega_0^2z=0
\]
reduces to
\(\ddot w+(\omega_0^2+\Omega_n^2)w=0\).
\end{proof}

For a spherical body whose spin axis defines latitude \(\lambda\),
\[
\Omega_n=\Omega\sin\lambda.
\]
Hence the magnitude vanishes at the equator and equals
\(\lvert\Omega\rvert\) at the poles. The sign depends on the
orientation of \(n,e_1,e_2\) and on the sign of the body's rotation;
only after these are declared is ``clockwise'' meaningful.

\section{Dissipation and energy}

Define the reduced mechanical energy
\[
E_{\rm loc}
=\frac m2\norm{\dot u}^2+\frac{m\omega_0^2}{2}\norm u^2.
\]

\begin{proposition}[Coriolis work and damping sign]
Every solution of \cref{eq:local-foucault} satisfies
\begin{equation}\label{eq:local-energy}
\dot E_{\rm loc}
=-2m\beta\norm{\dot u}^2.
\end{equation}
Thus the Coriolis term performs no work and
\[
P_{\rm diss}=2m\beta\norm{\dot u}^2\ge0.
\]
\end{proposition}

\begin{proof}
Dot \cref{eq:local-foucault} with \(m\dot u\). Since
\(\dot u\cdot J\dot u=0\), the Coriolis contribution vanishes.
\end{proof}

Equation \cref{eq:local-energy} does not imply that the full
body--pendulum--support energy is autonomous. Body rotation and
suspension control are externalized by the reduced model.

\section{Nested torus layer}

Let the active phases be
\[
\Theta(t)=\Theta_0+\omega t\pmod{2\pi},
\qquad \Theta\in\T^m .
\]
A finite circular hierarchy and local oscillation can be represented
by a continuous map \(F:\T^m\to\R^3\), with ideal trace
\[
Y(t)=F(\Theta(t)).
\]
This includes finite Fourier readouts
\[
Y(t)=\sum_{k\in\mathcal K}
c_ke^{\,i(k\cdot\omega)t},
\qquad \mathcal K\subset\mathbb Z^m\ \text{finite}.
\]

\begin{theorem}[Orbit closure of the phase flow]
Let
\[
\Lambda_\omega
=\{k\in\mathbb Z^m:k\cdot\omega=0\}.
\]
The closure of \(\Theta(t)\) is the translate by \(\Theta_0\) of the
closed subtorus whose annihilator is \(\Lambda_\omega\). In
particular:
\begin{enumerate}[label=(\roman*)]
  \item the flow is periodic iff some \(T>0\) satisfies
  \(T\omega\in2\pi\mathbb Z^m\);
  \item if the components of \(\omega\) are rationally independent,
  the phase orbit is dense in \(\T^m\);
  \item \(\cl\{Y(t)\}=F(\cl\{\Theta(t)\})\).
\end{enumerate}
\end{theorem}

\begin{proof}
The first two statements are the standard character description of a
one-parameter subgroup of a torus. The last follows because the phase
closure is compact and \(F\) is continuous.
\end{proof}

The map \(F\) can collapse dimensions or frequency amplitudes.
Therefore a full phase torus need not be identifiable from one
projected trace.

\section{Transformation audit}

\begin{longtable}{L{34mm}L{48mm}L{68mm}}
\toprule
Quantity or structure & Declared transformation & Status \\
\midrule
\endhead
\(\Omega^B\cdot n^B\) & common \(A\in\SO(3)\) applied to both vectors
& invariant scalar \\
Endpoint \(p^I\) & affine Euclidean frame change & covariant point,
not invariant components \\
Local energy \(E_{\rm loc}\) & static tangent-plane orthogonal basis
change & invariant \\
Orbit-closure topology & static observer-space homeomorphism &
preserved \\
Frequency set & static Euclidean isometry and time-origin shift &
preserved up to zero/trend conventions \\
Frequency amplitudes & camera projection & may be suppressed \\
Frequency set & arbitrary time-dependent rotating observer & not
invariant; modulation and shifts may occur \\
Signed precession & reversal of tangent-plane orientation & changes
sign \\
\bottomrule
\end{longtable}

\begin{warning}[No universal spectral invariant]
``Frequency support'' is not an invariant under every observer
choice. The admissible statement requires a fixed clock and a static
Euclidean change of observer, or a specified correction for the
observer motion.
\end{warning}

\section{Non-identifiability}

\begin{proposition}[One trace does not identify one hierarchy]
For a single linear camera projection and finite observation horizon,
the map from hierarchy parameters, local site state, and phase
amplitudes to the projected trace is not injective in general.
\end{proposition}

\begin{proof}
Any component in the kernel of the camera projection is invisible.
Furthermore, two latent Fourier channels whose projected coefficient
sum is equal produce the same projected signal when their frequencies
coincide. Finite sampled data additionally admit aliases separated by
integer multiples of the sampling angular frequency. Hence distinct
hidden records can have the same observation.
\end{proof}

Recoverability requires extra hypotheses: multiple noncoplanar camera
channels, multiple sites, known body attitudes, anti-alias sampling,
or external ephemeris constraints. The present document proves none
of these sufficient in a realistic network.

\section{Analytic sanity checks}

Using the conventional nominal terrestrial rotation magnitude
\(\Omega_\oplus=7.2921150\times10^{-5}\,\mathrm{s}^{-1}\), the
reduced-model precession period is
\[
T_F(\lambda)=\frac{2\pi}{\Omega_\oplus\lvert\sin\lambda\rvert}
\]
when \(\sin\lambda\ne0\). This gives:

\begin{center}
\begin{tabular}{ccc}
\toprule
Latitude & \(\Omega_n/\Omega_\oplus\) & \(T_F\) \\
\midrule
\(0^\circ\) & \(0\) & no reduced-model precession \\
\(45^\circ\) & \(1/\sqrt2\) & \(33.85\) h \\
\(90^\circ\) & \(1\) & \(23.93\) h \\
\bottomrule
\end{tabular}
\end{center}

These are formula checks, not predictions for a real installation.
The nominal rotation value does not encode polar motion or
length-of-day variations.

\section{Claim boundary}

The strict contributions are frame correction, typed observation
maps, a proved precession statement inside a declared reduced model,
and qualified transformation laws. This document does not claim:
\begin{itemize}
  \item a new derivation of the Foucault effect;
  \item a KAM theorem for the Solar System;
  \item a realistic pendulum on another celestial body;
  \item recovery of a celestial hierarchy from one trace;
  \item an empirical test or a new physical parameter.
\end{itemize}

\begingroup
\hbadness=10000
\begin{thebibliography}{9}

\bibitem{arnold}
V.~I. Arnold, \emph{Mathematical Methods of Classical Mechanics},
2nd ed., Springer, 1989.
\href{https://doi.org/10.1007/978-1-4757-2063-1}
{doi:10.1007/978-1-4757-2063-1}.

\bibitem{foucault-sagnac}
G. Rizzi and M.~L. Ruggiero, ``The Foucault pendulum and the Sagnac
effect,'' \emph{General Relativity and Gravitation} 35 (2003),
1743--1764.
\href{https://doi.org/10.1023/A:1021932917936}
{doi:10.1023/A:1021932917936}.

\bibitem{katok}
A. Katok and B. Hasselblatt,
\emph{Introduction to the Modern Theory of Dynamical Systems},
Cambridge University Press, 1995.

\end{thebibliography}
\endgroup

\end{document}
